\chapter{Factibilidad}
\label{cap:factibilidad}
El análisis de factibilidad examina si puedo desarrollar y evaluar Neuma con los
recursos y el tiempo disponibles. El problema y el alcance se presentan en los
capítulos~\ref{cap:introduccion} y~\ref{cap:objetivos}; aquí considero las
condiciones que permiten llevarlos adelante y los riesgos que pueden afectar
su cumplimiento. Distingo la previsión inicial del protocolo de la situación
conocida al redactar este capítulo, el 4 de noviembre de 2026, cuando ya existe
un prototipo implementado.

\section{Factibilidad técnica}
Dispongo de una notebook MSI Thin GF63 con procesador Intel Core i5-10500H,
16 GB de RAM DDR4 y una GPU NVIDIA RTX 3050 con 4 GB de memoria gráfica.
También dispongo de un Samsung S22 Ultra para las pruebas en un teléfono físico.
Ambos equipos eran de mi propiedad antes de iniciar el proyecto y no necesité
comprar equipamiento específico. La computadora proporciona el entorno de
desarrollo y construcción; el teléfono permite abordar la captura, reproducción
e interacción bajo condiciones de uso reales.

La aplicación utiliza las tecnologías descritas en el capítulo~\ref{cap:herramientas}.
La existencia de una candidata compilada y de verificaciones automatizadas de
reglas y componentes respalda la posibilidad de construir e integrar la solución.
Los registros de la candidata 0.12.20 distinguen esas comprobaciones en computadora
de las verificaciones Android y acústicas todavía pendientes. Los resultados y
condiciones de cada prueba se reúnen en el capítulo~\ref{cap:pruebas}.

La principal incertidumbre técnica se concentra en el comportamiento del audio
con voz real, ruido y diferencias entre dispositivos, y en el ajuste de los
parámetros de adaptación. Mi experiencia previa en desarrollo Android y audio
en tiempo real es limitada; la construcción incremental, la revisión de propuestas
y las pruebas permiten abordar esa complejidad por partes. La disponibilidad
del S22 Ultra facilita una primera evaluación física, aunque la cobertura de
otros equipos dependerá de los dispositivos a los que pueda acceder.

El protocolo inicial propone una arquitectura con C++/Oboe y procesamiento
híbrido DSP/ML. En el desarrollo adopté primero la línea base en Kotlin con YIN,
tratando las alternativas como ampliaciones condicionadas a mediciones, según
lo expuesto en el alcance. La factibilidad de la entrega depende de comprobar
la suficiencia de esa línea base y acordar con el director el tratamiento de
las diferencias respecto del protocolo. La incorporación de nuevas tecnologías
deberá responder a una limitación identificada y considerar el tiempo de
integración y verificación que requiera.

\section{Factibilidad operativa}
El uso previsto complementa otras formas de estudio musical y se apoya en un
teléfono Android, sus medios de reproducción y, para los ejercicios vocales,
su micrófono. La ejecución local y el perfil almacenado en el dispositivo
permiten plantear la práctica sin una cuenta remota ni conexión permanente.
Los ejercicios de voz requieren condiciones acústicas que permitan obtener
una señal utilizable; la comprensión de las consignas y devoluciones también
debe examinarse con participantes.

Cuento con personas a quienes puedo proponer las pruebas. Esto proporciona una
vía de convocatoria, cuya concreción dependerá de su interés, disponibilidad y
aceptación de las condiciones. Las modalidades autónoma mediante APK y presencial
acompañada permiten contemplar distintas disponibilidades. Las tareas, registros
y criterios se desarrollan en la sección~\ref{sec:metodo-evaluacion}.

En la modalidad autónoma, la instalación y actualización de la APK forman parte
de las condiciones de acceso al prototipo y pueden requerir orientación.
Antes de convocar las pruebas deberé disponer de una versión suficientemente
estable y de instrucciones de instalación y uso. El relevamiento preparado y
las pruebas previstas permitirán revisar las necesidades y dificultades; su
preparación todavía no constituye evidencia de aceptación o utilidad del producto.

\section{Factibilidad temporal}
Actualmente dedico aproximadamente 12 horas semanales al proyecto. Mantengo
fines de febrero de 2027 como fecha prevista de finalización y marzo de 2027
como mes provisional de defensa. El protocolo fija el cierre el 27 de febrero;
entre la fecha de este análisis y ese día hay unas 16 semanas, equivalentes a
aproximadamente 200 horas de dedicación potencial si sostengo el ritmo declarado.
Esta proyección no descuenta interrupciones y comprende el conjunto de tareas:
desarrollo, pruebas, análisis, escritura y revisiones.

El plan original contempla 400 horas: 30 previas estimadas, 335 previstas y
35 de reserva, con una capacidad declarada de 10 a 20 horas semanales. La
dedicación actual se encuentra dentro de ese intervalo, pero exige revisar
la distribución de las semanas que habían sido planificadas con una carga mayor.
Las 400 horas son una estimación de partida; el avance del prototipo no permite
deducir por sí solo cuántas horas restan ni qué parte de la reserva sigue disponible.

La continuidad hasta el cierre requiere priorizar la estabilización del alcance
evaluable, las mediciones físicas, las pruebas de uso y la terminación del informe.
Las ampliaciones competirán por esa misma disponibilidad y deberán revisarse
según su aporte al cumplimiento de los objetivos. El capítulo~\ref{cap:planificacion}
desarrollará las tareas, dependencias y diferencias entre el plan inicial y la
ejecución; la compatibilidad temporal deberá actualizarse con esa revisión y
con los acuerdos del director.

\section{Factibilidad económica}
La infraestructura principal está disponible y no requirió compras específicas
para el proyecto. Esto reduce la necesidad de un desembolso inicial en equipos,
aunque su uso y mi dedicación siguen siendo recursos empleados en el trabajo.
El prototipo funciona en el dispositivo y no requiere contratar un servidor para
ejecutar los ejercicios y conservar el perfil local.

Al momento de este análisis pago USD 20 mensuales por Codex y USD 30 mensuales
por Grok, que suman USD 50 mensuales en suscripciones. Ese monto identifica el
gasto mensual en servicios de apoyo. Para determinar un costo acumulado atribuible
a Neuma sería necesario precisar los períodos y la proporción de uso dedicada
al proyecto; por ello, no lo presento como costo total del Trabajo Final. La
continuidad de los servicios de apoyo implica sostener ese gasto o revisar las
herramientas utilizadas si cambia su disponibilidad.

El análisis se limita a los recursos disponibles y los gastos identificados.
No asigno un valor monetario a mis horas ni estimo amortización de los equipos.
La entrega gratuita prevista para quienes utilicen Neuma es una decisión de
acceso al producto, compatible con que su desarrollo demande tiempo y servicios
de apoyo. La publicación abierta requerirá completar la revisión de licencias
de código y recursos.

\section{Riesgos}
La tabla~\ref{tab:riesgos-factibilidad} reúne riesgos que pueden afectar el
cierre y las medidas previstas para tratarlos. Su identificación es cualitativa,
a partir del protocolo, el estado del prototipo y los recursos disponibles;
no se dispone de datos para asignar probabilidades numéricas. Las respuestas
indicadas orientan el trabajo pendiente y deberán actualizarse con las pruebas
y revisiones.

\begingroup
\small
\singlespacing
\setlength{\tabcolsep}{4pt}
\setlength{\LTcapwidth}{\textwidth}
\begin{longtable}{p{0.22\textwidth}p{0.29\textwidth}p{0.39\textwidth}}
\caption{Riesgos de factibilidad y medidas previstas.}
\label{tab:riesgos-factibilidad}\\
\toprule
Riesgo & Efecto sobre la entrega & Medida prevista \\
\midrule
\endfirsthead
\toprule
Riesgo & Efecto sobre la entrega & Medida prevista \\
\midrule
\endhead
\bottomrule
\endfoot
Audio y diversidad de equipos & Estimaciones inestables o respuesta tardía que dificulten los ejercicios vocales. & Medir con referencias y voz en condiciones registradas; corregir la línea base y delimitar las condiciones de funcionamiento. \\
Calibración del DDA & Ajustes poco adecuados o progresión inconsistente. & Revisar entradas, reglas y decisiones; comprobar límites y ajustar parámetros a partir de los registros. \\
Errores en propuestas o código asistido & Fallos de integración o explicaciones incorrectas. & Mantener especificaciones, revisión humana, pruebas y contraste de fuentes; conservar versiones recuperables. \\
Ampliación del alcance & Desplazamiento de mediciones, escritura y revisión. & Priorizar el alcance evaluable y revisar nuevas funciones con relación al tiempo disponible. \\
Convocatoria y coordinación & Menos pruebas de uso o demoras para realizarlas. & Convocar con anticipación, ofrecer ambas modalidades y acordar con el director el procedimiento y las revisiones. \\
Servicios de apoyo y publicación & Interrupciones del trabajo o demoras por recursos sin licencia aclarada. & Conservar código y documentación local; revisar herramientas si cambia su acceso y auditar licencias antes de publicar. \\
\end{longtable}
\endgroup

\section{Síntesis de factibilidad}
La infraestructura disponible, la base de software implementada y el acceso a
posibles participantes permiten continuar con la construcción y evaluación del
prototipo. Mi dedicación y los gastos de apoyo identificados proporcionan un
marco concreto para organizar esa continuidad. La viabilidad del cierre previsto
queda condicionada a estabilizar el alcance, completar la evaluación física y
de uso y reservar tiempo para el informe y su revisión con el director.
